% =====================================================================
%  alcance.tex  --  Documento de Alcance de LEIP
%  ARCHIVO UNICO: no necesita ningun otro archivo.
%  Compilar en Overleaf con pdfLaTeX (dos pasadas por el indice).
% =====================================================================
\documentclass[11pt,a4paper]{article}

% ---------------------------------------------------------------------
%  DATOS DEL PROYECTO  (tocá solo esto para actualizar el documento)
% ---------------------------------------------------------------------
\newcommand{\proyectoNombre}{LEIP}
\newcommand{\proyectoNombreLargo}{LatAm Economic Intelligence Platform}
\newcommand{\proyectoUniversidad}{Universidad Cat\'olica del Uruguay}
\newcommand{\proyectoFacultad}{Facultad de Ingenier\'ia y Tecnolog\'ias}
\newcommand{\proyectoCarrera}{Ingenier\'ia en Inform\'atica}
\newcommand{\docVersion}{1.0}
\newcommand{\docTitulo}{Alcance}
\newcommand{\docFecha}{Agosto 2026}
\newcommand{\modMacro}{Macro Expectations Intelligence}
\newcommand{\modPrecios}{Prices \& Inflation Intelligence}
\newcommand{\modActividad}{Real Activity Intelligence}

% ---------------------------------------------------------------------
%  PREAMBULO  (estilos y comandos; normalmente no hace falta tocarlo)
% ---------------------------------------------------------------------
\usepackage[utf8]{inputenc}
\usepackage[T1]{fontenc}
\IfFileExists{lmodern.sty}{\usepackage{lmodern}\usepackage{microtype}}{}

\IfFileExists{spanish.ldf}{%
  \usepackage[spanish,es-tabla,es-noquoting]{babel}%
}{%
  \usepackage[english]{babel}%
  \renewcommand{\contentsname}{\'Indice}%
  \renewcommand{\listfigurename}{\'Indice de figuras}%
  \renewcommand{\listtablename}{\'Indice de tablas}%
  \renewcommand{\tablename}{Tabla}%
  \renewcommand{\figurename}{Figura}%
  \renewcommand{\refname}{Referencias}%
}

\usepackage[a4paper,margin=2.5cm,headheight=15pt]{geometry}
\usepackage{fancyhdr}
\usepackage{parskip}
\usepackage{booktabs}
\usepackage{tabularx}
\usepackage{longtable}
\usepackage{array}
\usepackage{enumitem}
\setlist{noitemsep,topsep=4pt}
\usepackage{graphicx}
\usepackage{xcolor}
\usepackage{tcolorbox}
\tcbuselibrary{breakable}
\usepackage[hidelinks]{hyperref}
\usepackage{url}
\urlstyle{same}

\definecolor{leipAzul}{HTML}{1F4E79}
\definecolor{leipGris}{HTML}{5A5A5A}
\definecolor{leipAmbar}{HTML}{B8860B}
\definecolor{leipAmbarBg}{HTML}{FFF8E1}

\pagestyle{fancy}
\fancyhf{}
\fancyhead[L]{\small\color{leipGris}\proyectoNombre{} --- \docTitulo}
\fancyhead[R]{\small\color{leipGris} v\docVersion}
\fancyfoot[C]{\small\thepage}
\renewcommand{\headrulewidth}{0.4pt}

\usepackage{titlesec}
\titleformat{\section}{\Large\bfseries\color{leipAzul}}{\thesection}{0.6em}{}
\titleformat{\subsection}{\large\bfseries\color{leipAzul}}{\thesubsection}{0.6em}{}
\titleformat{\subsubsection}{\normalsize\bfseries\color{leipGris}}{\thesubsubsection}{0.6em}{}

% --- Portada -----------------------------------------------------------
\newcommand{\portada}[1]{%
  \begin{titlepage}
  \centering
  \vspace*{2cm}
  {\Large \proyectoUniversidad\par}
  {\large \proyectoFacultad\par}
  {\large \proyectoCarrera\par}
  \vspace{3cm}
  {\Huge\bfseries\color{leipAzul} #1\par}
  \vspace{0.8cm}
  {\Large \proyectoNombre{} --- \proyectoNombreLargo\par}
  \vspace{2cm}
  {\large Versi\'on \docVersion\par}
  {\large \docFecha\par}
  \vfill
  {\large\bfseries Integrantes\par}
  \vspace{0.3cm}
  \begin{tabular}{c}
    Sebastian Forische \\
    Armando Hernandez de la Vega \\
    Luana Acu\~na \\
    Bruno Sebastian Olivera Viera Gomez \\
  \end{tabular}
  \vspace{1.5cm}
  \end{titlepage}
}

% --- Historial de revisiones -------------------------------------------
\newenvironment{historial}{%
  \section*{Historial de revisiones}
  \addcontentsline{toc}{section}{Historial de revisiones}
  \begin{tabular}{@{}l l p{0.52\textwidth} l@{}}
  \toprule
  \textbf{Versi\'on} & \textbf{Fecha} & \textbf{Descripci\'on} & \textbf{Autor} \\
  \midrule
}{%
  \bottomrule
  \end{tabular}
  \bigskip
}

% --- Requerimientos numerados ------------------------------------------
% Uso:  \rf{etiqueta}{texto del requerimiento}
%       \rnf{etiqueta}{texto del requerimiento}
% Después se referencian con \ref{rf:etiqueta} y \ref{rnf:etiqueta}
\newcounter{rfcounter}
\newcounter{rnfcounter}
\newcommand{\dosdig}[1]{\ifnum\value{#1}<10 0\fi\arabic{#1}}

\newcommand{\rf}[2]{%
  \refstepcounter{rfcounter}\label{rf:#1}%
  \par\smallskip\noindent
  \begin{minipage}{\textwidth}
  \hangindent=1.9cm \hangafter=1
  \makebox[1.7cm][l]{\textbf{\color{leipAzul}RF-\dosdig{rfcounter}}}#2
  \end{minipage}\par
}
\newcommand{\rnf}[2]{%
  \refstepcounter{rnfcounter}\label{rnf:#1}%
  \par\smallskip\noindent
  \begin{minipage}{\textwidth}
  \hangindent=2.1cm \hangafter=1
  \makebox[1.9cm][l]{\textbf{\color{leipAzul}RNF-\dosdig{rnfcounter}}}#2
  \end{minipage}\par
}
\renewcommand{\therfcounter}{RF-\dosdig{rfcounter}}
\renewcommand{\thernfcounter}{RNF-\dosdig{rnfcounter}}

% --- Figura con fallback si la imagen no existe ------------------------
\newcommand{\figuraLEIP}[4]{%
  \begin{figure}[htbp]
  \centering
  \IfFileExists{#1}{\includegraphics[width=#4\textwidth]{#1}}{%
    \fbox{\parbox[c][4cm][c]{0.8\textwidth}{\centering\color{leipAmbar}%
      \textbf{[Figura pendiente]}\\[4pt]\small\texttt{#1}}}%
  }
  \caption{#2}
  \label{#3}
  \end{figure}
}

\begin{document}

\portada{Alcance}

\tableofcontents
\bigskip

\begin{historial}
1.0 & 26/06/2026 & Contenido inicial derivado del anteproyecto & Equipo \\
\end{historial}
\clearpage

% =====================================================================
\section{Introducci\'on}

\subsection{Prop\'osito}

Este documento define el alcance de \proyectoNombre{} (\proyectoNombreLargo):
qu\'e construye el proyecto y qu\'e queda expl\'icitamente afuera. Enumera los
requerimientos funcionales y no funcionales del sistema, identifica usuarios e
interesados, y establece las restricciones que condicionan el desarrollo.

Mientras que el Documento de Visi\'on responde \emph{por qu\'e} y \emph{para
qui\'en} existe el producto, este documento responde \emph{qu\'e} debe hacer
exactamente. Es el documento de referencia para acordar el alcance con la
contraparte y la base sobre la que se derivan los casos de uso, los casos de
prueba y las decisiones de arquitectura.

\subsection{Relaci\'on con los dem\'as documentos}

\begin{tabularx}{\textwidth}{l X}
\toprule
\textbf{Documento} & \textbf{Relaci\'on} \\
\midrule
Visi\'on & Justifica el problema, el mercado y el posicionamiento. Este documento no repite ese contenido. \\
Casos de Uso & Detalla la interacci\'on de cada actor con el sistema. Cada caso de uso se deriva de uno o m\'as requerimientos funcionales de este documento. \\
Casos de Prueba & Define c\'omo se verifica cada caso de uso. \\
Notas de Arquitectura y Dise\~no & Justifica cada decisi\'on t\'ecnica contra los requerimientos no funcionales de la secci\'on \ref{sec:rnf}. \\
Lista de Riesgos & Registra los riesgos asociados a las asunciones y restricciones de este documento. \\
\bottomrule
\end{tabularx}

\subsection{Convenci\'on de identificadores}

Los requerimientos se identifican como \textbf{RF-nn} (funcionales) y
\textbf{RNF-nn} (no funcionales). Estos identificadores son estables: una vez
asignados no se reutilizan ni se renumeran, aunque el requerimiento se elimine.
Los casos de uso, casos de prueba y decisiones de arquitectura referencian estos
identificadores.

% =====================================================================
\section{Propuesta}

\proyectoNombre{} es una plataforma SaaS de inteligencia econ\'omica para
Am\'erica Latina que transforma informaci\'on oficial dispersa y heterog\'enea
en una capa comparable, trazable y actualizada.

La plataforma obtiene datos desde fuentes oficiales mediante procesos de ingesta
que identifican la fuente, extraen informaci\'on desde distintos formatos,
validan su contenido y normalizan variables, unidades, fechas, frecuencias y
metodolog\'ias. Los datos procesados se organizan por pa\'is, variable,
per\'iodo, unidad, fuente, fecha de publicaci\'on y versi\'on. Sobre esa base se
construye una capa de inteligencia que calcula m\'etricas determin\'isticas y
cruces entre m\'odulos.

El usuario accede mediante una aplicaci\'on web con \emph{dashboards},
comparadores, fichas por pa\'is, radar, alertas, \emph{briefings} y
exportaciones auditables. Las organizaciones pueden adem\'as consumir una API
institucional autenticada.

% =====================================================================
\section{Usuarios e interesados}

\subsection{Interesados}

\subsection{Tipos de usuario}

Los siguientes tipos de usuario se derivan del mercado objetivo definido en el
Documento de Visi\'on. Cada uno se corresponde con actores de la Especificaci\'on
de Casos de Uso.

\begin{tabularx}{\textwidth}{l X l}
\toprule
\textbf{Tipo de usuario} & \textbf{Relaci\'on con el sistema} & \textbf{Plan t\'ipico} \\
\midrule
Usuario Free & Consulta el radar resumido y un conjunto limitado de indicadores. Es la puerta de entrada al producto. & Free \\
\addlinespace
Analista individual & Consulta, compara y exporta datos y m\'etricas de los m\'odulos contratados. & M\'odulo individual / Completo \\
\addlinespace
Usuario de organizaci\'on & Igual que el analista, dentro de una cuenta corporativa con administraci\'on centralizada. & Team / Enterprise \\
\addlinespace
Administrador de organizaci\'on & Gestiona usuarios, permisos y suscripci\'on de su propia organizaci\'on. & Team / Enterprise \\
\addlinespace
Sistema consumidor de API & Sistema externo de una instituci\'on que consume datos y m\'etricas de forma programada. & Team / Enterprise \\
\addlinespace
Administrador de \proyectoNombre{} & Opera la plataforma: gestiona cuentas, planes, m\'odulos, fuentes y cuarentena. & --- \\
\addlinespace
Proveedor de pagos & Sistema externo que notifica cambios de estado de suscripci\'on mediante \emph{webhooks}. & --- \\
\addlinespace
Programador de tareas & Actor no humano que dispara los procesos de ingesta y c\'alculo seg\'un la frecuencia de cada fuente. & --- \\
\bottomrule
\end{tabularx}

% =====================================================================
\section{Contexto del producto}

\proyectoNombre{} opera como sistema independiente que se relaciona con cuatro
tipos de entidades externas:

\begin{itemize}
  \item \textbf{Organismos oficiales emisores:} BCB (Brasil), Banxico (M\'exico)
        y BCU (Uruguay). Publican la informaci\'on de origen. La relaci\'on es
        unidireccional y sin acuerdo de servicio: el sistema consume lo que
        el organismo publica, en el formato en que lo publica.
  \item \textbf{Proveedor externo de cobros:} aloja el \emph{checkout}, procesa
        los pagos y notifica los cambios de estado de suscripci\'on.
  \item \textbf{Sistemas consumidores de la API institucional:} sistemas propios
        de organizaciones cliente.
  \item \textbf{Fuentes privadas de clientes institucionales:} datos aportados
        por el propio cliente, procesados en entorno aislado.
\end{itemize}

\figuraLEIP{img/contexto.png}{Diagrama de contexto de \proyectoNombre{}}{fig:contexto}{0.9}

% =====================================================================
\section{Capacidades del MVP}

\begin{itemize}
  \item \textbf{Pa\'ises:} Brasil, M\'exico y Uruguay.
  \item \textbf{M\'odulos:} \modMacro{}, \modPrecios{} y \modActividad{}.
  \item \textbf{M\'etricas determin\'isticas:} cambios, momentum, divergencias
        regionales, percentiles hist\'oricos, distancia frente a metas,
        sorpresas y persistencia.
  \item \textbf{Cruces:} inflaci\'on esperada frente a observada, y crecimiento
        esperado frente a actividad efectiva.
  \item \textbf{Aplicaci\'on web:} \emph{dashboards}, comparadores, fichas por
        pa\'is, radar, alertas, \emph{briefings} y exportaciones auditables.
  \item \textbf{Control de acceso:} autenticaci\'on, roles, permisos y
        \emph{entitlements} por plan, m\'odulo, pa\'is y funcionalidad.
  \item \textbf{API institucional} autenticada, con cuotas y l\'imites por plan.
\end{itemize}

% =====================================================================
\section{Requerimientos funcionales}
\label{sec:rf}

% ---------------------------------------------------------------------
\subsection{Sistema de datos}

\subsubsection{Ingesta de fuentes oficiales}

\rf{ingesta-obtener}{El sistema debe obtener informaci\'on econ\'omica oficial
para los m\'odulos incluidos en el MVP, desde las fuentes definidas para cada
pa\'is.}

\rf{ingesta-programada}{El sistema debe ejecutar actualizaciones autom\'aticas
seg\'un la frecuencia de publicaci\'on de cada fuente.}

\rf{ingesta-registro}{El sistema debe registrar, para cada dato o archivo
incorporado: fuente, fecha de publicaci\'on, fecha de ingesta, formato,
versi\'on y \emph{checksum}.}

\subsubsection{Normalizaci\'on de datos}

\rf{norm-modelo}{El sistema debe estandarizar pa\'ises, variables, unidades,
monedas, fechas, frecuencias y horizontes mediante un modelo com\'un.}

\rf{norm-formatos}{El sistema debe transformar informaci\'on publicada en API,
Excel, CSV, PDF, HTML y JSON hacia el modelo com\'un.}

\rf{norm-historico}{El sistema debe reconstruir series hist\'oricas, conservar
las revisiones sucesivas de cada dato y mantener el estado de la informaci\'on
disponible en cada momento.}

\subsubsection{Control de calidad y trazabilidad}

\rf{qa-deteccion}{El sistema debe detectar duplicados, faltantes,
inconsistencias y cambios de formato antes de publicar informaci\'on.}

\rf{qa-cuarentena}{El sistema debe enviar los registros dudosos a cuarentena,
conservar los originales inmutables en la capa Bronze y registrar ejecuciones,
errores y correcciones.}

% ---------------------------------------------------------------------
\subsection{Capa de inteligencia}

\subsubsection{M\'etricas}

\rf{met-calculo}{El sistema debe calcular cambios, momentum, divergencias
regionales, percentiles hist\'oricos, sorpresas y persistencia sobre las series
normalizadas.}

\rf{met-config}{El sistema debe permitir configurar ventanas, par\'ametros y
metodolog\'ias de c\'alculo seg\'un cada m\'odulo.}

\rf{met-metas-cruces}{El sistema debe calcular la distancia frente a metas y
resolver los cruces entre expectativas, inflaci\'on observada y actividad
efectiva.}

\subsubsection{Comparaci\'on regional}

\rf{comp-paises}{El sistema debe permitir comparar pa\'ises, variables,
per\'iodos y m\'odulos econ\'omicos de Am\'erica Latina bajo criterios
homog\'eneos.}

\rf{comp-rankings}{El sistema debe permitir identificar revisiones, divergencias
y \emph{rankings} regionales relevantes.}

\rf{comp-advertencia}{El sistema debe advertir al usuario cuando un cruce no
pueda calcularse, distinguiendo si la causa es la falta de datos o la ausencia
de los m\'odulos requeridos en el plan contratado.}

\subsubsection{Se\~nales anal\'iticas}

\rf{senal-alertas}{El sistema debe generar alertas e informes descriptivos sobre
cambios relevantes previamente calculados y validados.}

\rf{senal-clasificacion}{El sistema debe clasificar los movimientos seg\'un su
direcci\'on, intensidad y relevancia para el an\'alisis.}

\rf{senal-resumen}{El sistema debe resumir qu\'e se movi\'o en la regi\'on sin
que el usuario deba descargar, limpiar o reconstruir manualmente los datos. El
resumen debe derivarse exclusivamente de m\'etricas ya calculadas y validadas.}

% ---------------------------------------------------------------------
\subsection{Aplicaci\'on web}

\subsubsection{Autenticaci\'on y usuarios}

\rf{auth-registro}{El sistema debe permitir a los usuarios registrarse e iniciar
sesi\'on con sus credenciales.}

\rf{auth-recuperacion}{El sistema debe permitir al usuario recuperar el acceso a
su cuenta.}

\rf{auth-entitlements}{El sistema debe aplicar roles, permisos y
\emph{entitlements} seg\'un el plan, los m\'odulos, los pa\'ises y las
funcionalidades contratadas.}

\rf{auth-admin}{El sistema debe permitir al administrador gestionar usuarios,
organizaciones, planes, m\'odulos, permisos y suscripciones.}

\subsubsection{Pagos y suscripciones}

\rf{pago-checkout}{El sistema debe delegar el pago al \emph{checkout} de un
proveedor externo. \proyectoNombre{} no almacena ni procesa directamente datos
de tarjetas.}

\rf{pago-webhook}{El sistema debe recibir y procesar las notificaciones del
proveedor de pagos cuando una suscripci\'on sea aprobada, renovada, rechazada,
cancelada o reembolsada.}

\rf{pago-datos}{El sistema debe conservar \'unicamente los identificadores
externos, el estado de la suscripci\'on y los datos necesarios para actualizar
m\'odulos, planes y permisos.}

\rf{pago-cambio-plan}{El sistema debe reflejar los cambios de plan en los
\emph{entitlements} de la cuenta, habilitando o revocando el acceso a m\'odulos y
funcionalidades seg\'un corresponda.}

\subsubsection{Dashboard de an\'alisis}

\rf{dash-seleccion}{El sistema debe permitir al usuario seleccionar pa\'is,
m\'odulo, variable, fuente y per\'iodo hist\'orico.}

\rf{dash-visualizacion}{El sistema debe permitir visualizar gr\'aficos, tablas,
m\'etricas, \emph{rankings} y cruces entre m\'odulos desde una misma interfaz.}

\rf{dash-navegacion}{El sistema debe permitir navegar entre pa\'ises, variables,
m\'odulos y a\~nos de forma simple y r\'apida.}

\rf{dash-historico}{El sistema debe permitir consultar hist\'oricos y revisiones
para analizar cambios y tendencias en el tiempo.}

\rf{dash-procedencia}{El sistema debe mostrar, junto a cada dato y m\'etrica
visualizada, su fuente, fecha, versi\'on y nota metodol\'ogica.}

\subsubsection{Radar de expectativas}

\rf{radar-acceso}{El sistema debe publicar un radar peri\'odico con los
principales cambios econ\'omicos de Am\'erica Latina.}

\rf{radar-free}{El sistema debe mostrar una versi\'on resumida del radar a los
usuarios del plan Free, sin requerir contrataci\'on.}

\rf{radar-pago}{El sistema debe mostrar la versi\'on completa del radar, las
alertas y los \emph{briefings} a los usuarios con plan de pago vigente.}

\subsubsection{Exportaci\'on y uso de resultados}

\rf{exp-resultados}{El sistema debe permitir a los usuarios de planes
profesionales exportar resultados transformados para an\'alisis, estudio o
reportes personales.}

\rf{exp-formatos}{El sistema debe permitir descargar informaci\'on en formatos
de uso corriente, como Excel y PDF.}

\rf{exp-trazabilidad}{El sistema debe incluir en cada exportaci\'on la fuente,
fecha, versi\'on y nota metodol\'ogica de los datos y m\'etricas exportados.}

\rf{exp-consultas}{El sistema debe permitir guardar consultas frecuentes para
facilitar el seguimiento recurrente de expectativas.}

\subsubsection{API institucional y fuentes privadas}

\rf{api-consulta}{El sistema debe permitir a las instituciones consultar datos y
m\'etricas mediante una API autenticada, con cuotas, l\'imites y permisos seg\'un
el plan contratado.}

\rf{api-privadas}{El sistema debe permitir a las instituciones integrar fuentes
privadas dentro de un entorno aislado, sin mezclarlas con la informaci\'on de
otros clientes.}

% =====================================================================
\section{Requerimientos no funcionales}
\label{sec:rnf}

Los siguientes requerimientos condicionan las decisiones de arquitectura. Cada
decisi\'on documentada en Notas de Arquitectura y Dise\~no debe indicar a
cu\'al de estos responde.

\subsection{Trazabilidad y auditor\'ia}

\rnf{traza-completa}{Cada dato, m\'etrica y exportaci\'on debe conservar su
fuente, fecha, versi\'on y criterio de procesamiento.}

\rnf{traza-inmutable}{Los archivos y respuestas originales de las fuentes deben
conservarse inmutables en la capa Bronze.}

\rnf{traza-auditoria}{El sistema debe mantener registros de auditor\'ia.}

\subsection{Seguridad}

\rnf{seg-superficie}{La plataforma debe exponer \'unicamente el frontend y los
\emph{endpoints} necesarios.}

\rnf{seg-sesiones}{El sistema debe utilizar sesiones o \emph{tokens} seguros y
validar las solicitudes recibidas.}

\rnf{seg-secretos}{La gesti\'on de secretos debe realizarse fuera del c\'odigo
fuente.}

\rnf{seg-permisos}{El sistema debe aplicar roles y permisos por usuario,
organizaci\'on, plan, m\'odulo y pa\'is.}

\rnf{seg-pci}{\proyectoNombre{} no debe almacenar ni procesar directamente datos
de tarjetas.}

\subsection{Aislamiento}

\rnf{multi-aislamiento}{El sistema debe proveer aislamiento
multi-\emph{tenant} para las fuentes privadas.}

\subsection{Cumplimiento}

\rnf{cump-transformado}{El producto debe entregar informaci\'on transformada
---m\'etricas, comparaciones y valor agregado--- y no una copia del dato crudo.}

\rnf{cump-atribucion}{Cada visualizaci\'on, exportaci\'on o respuesta de API
debe incluir fuente, fecha, versi\'on y nota metodol\'ogica.}

\section{Fuera de alcance}

Los siguientes elementos quedan expl\'icitamente excluidos del alcance del MVP:

\begin{itemize}
  \item \textbf{Chile.} No se incorpora como pa\'is navegable, descargable ni
        comercializable hasta obtener permiso expl\'icito del BCCh.
  \item \textbf{M\'odulos de riesgo fiscal, sector externo y condiciones
        monetarias y financieras.} Forman parte de la visi\'on de producto pero
        no del MVP.
  \item \textbf{Activos de mercado.} Quedan para una fase posterior.
  \item \textbf{Recomendaciones personalizadas de inversi\'on.} \proyectoNombre{}
        entrega informaci\'on agregada y descriptiva. Esta exclusi\'on es
        deliberada y es lo que mantiene bajo el riesgo de asesor\'ia financiera.
  \item \textbf{Procesamiento directo de medios de pago.} Se delega
        \'integramente al proveedor externo.
\end{itemize}

% =====================================================================
\section{Precedencia y prioridad}

\section{Asunciones y dependencias}

\section{Restricciones}
\label{sec:restricciones}

\subsection{Restricciones regulatorias}

El riesgo relevante del proyecto es la redistribuci\'on y el uso comercial de
informaci\'on publicada por organismos oficiales. Los requerimientos
\ref{rnf:cump-transformado} y \ref{rnf:cump-atribucion} son la mitigaci\'on
principal de este riesgo.

\begin{tabularx}{\textwidth}{l X}
\toprule
\textbf{Pa\'is} & \textbf{Condici\'on} \\
\midrule
Brasil &
Parte de un marco de datos abiertos. A\'un as\'i, no debe redistribuirse la base
cruda como copia del Banco Central, sino trabajar sobre datos transformados con
atribuci\'on clara de la fuente. \\
\addlinespace
M\'exico &
El marco general es favorable para usar informaci\'on oficial, pero queda
pendiente confirmar las condiciones espec\'ificas. \\
\addlinespace
Uruguay &
El BCU tiene un marco de datos abiertos, pero hay que confirmar que esa licencia
cubre espec\'ificamente las encuestas de expectativas. \\
\addlinespace
Chile &
Caso m\'as restrictivo. El BCCh limita la comercializaci\'on o redistribuci\'on
de la informaci\'on, y transformar el dato no elimina ese riesgo. Fuera del
alcance del MVP. \\
\bottomrule
\end{tabularx}

\subsection{Restricciones t\'ecnicas}

\begin{itemize}
  \item Los cruces entre m\'odulos s\'olo pueden calcularse cuando la cuenta
        tiene acceso a todos los m\'odulos involucrados (\ref{rf:comp-advertencia}).
  \item Los m\'odulos individuales se contratan hasta un m\'aximo de tres; por
        encima de ese l\'imite corresponde el plan Completo.
\end{itemize}

% =====================================================================
\section{Criterio de aceptaci\'on}


\end{document}